% GENERATED FILE. Edit canonical lessons, then run npm run generate:books.
% canonical-source-hash: fnv1a64:c027fbcf2398b69c
% canonical-lessons: SA-C135-siksate, SA-C135-pathayati, SA-C135-vismarati, SA-C135-arohati, SA-C135-avarohati

\chapter{Learns, Teaches, Forgets, Climbs, Climbs down}
\label{ch:sa-learns-teaches-forgets-climbs-climbs-down}
\hlchaptermodality{135}

\begin{chapteropening}
\textbf{By the end of this chapter:} \emph{I can say learns, teaches, forgets, climbs and climbs down.}

Complete the last lesson of chapter 135: I can say learns, teaches, forgets, climbs and climbs down.
\end{chapteropening}

\section[\texorpdfstring{\sk{शिक्षते} (śikṣate)}{शिक्षते (śikṣate)}]{\sk{शिक्षते} (śikṣate) — learns}

\label{lesson:SA-C135-siksate}



\begin{quote}
Before the new one: say the Sanskrit for sends, then the Sanskrit for searches.
\end{quote}

\subsection*{You'll want to know: \sk{शिक्षते}}
\textbf{\sk{शिक्षते}} — \emph{śikṣate} — \textquotedblleft{}learns\textquotedblright{}.

\begin{cousinweb}
\emph{śikṣate} comes from the desiderative of the root √śak \textquotedblleft{}be able\textquotedblright{} (MW): \textquotedblleft{}wishes to be able\textquotedblright{}, and so \textquotedblleft{}learns\textquotedblright{}. It lives on as Hindi \emph{sīkhnā} \textquotedblleft{}to learn\textquotedblright{} (CDIAL 12430).
\end{cousinweb}

\begin{grammarlens}[title={What you've built}]
One more word for this chapter.
\end{grammarlens}

\subsection*{Your turn}
\begin{itemize}
\raggedright
  \item \emph{Say it:} \emph{śikṣate}
  \item \emph{Say it:} \emph{śikṣate}, once more
  \item \emph{Say it:} say \emph{śikṣate}
  \item \emph{Recall:} say the Sanskrit for sends, then the Sanskrit for searches, then say \emph{śikṣate} again
\end{itemize}

\begin{tcolorbox}[breakable,colback=teal!4,colframe=teal!35!black,title={Before you move on}]
What does \sk{शिक्षते} mean? (\textquotedblleft{}Learns\textquotedblright{}.) Say it once more.
\end{tcolorbox}
\section[\texorpdfstring{\sk{पाठयति} (pāṭhayati)}{पाठयति (pāṭhayati)}]{\sk{पाठयति} (pāṭhayati) — teaches}

\label{lesson:SA-C135-pathayati}



\begin{quote}
Before the new one: say the Sanskrit for searches, then the Sanskrit for learns.
\end{quote}

\subsection*{You'll want to know: \sk{पाठयति}}
\textbf{\sk{पाठयति}} — \emph{pāṭhayati} — \textquotedblleft{}teaches\textquotedblright{}.

\begin{cousinweb}
\emph{pāṭhayati} is the causative of the root √paṭh \textquotedblleft{}read aloud, recite\textquotedblright{} (CDIAL 8037): \textquotedblleft{}makes someone recite\textquotedblright{}, and so \textquotedblleft{}teaches\textquotedblright{}. It is the causing form of \emph{paṭhati} \textquotedblleft{}reads\textquotedblright{}, and Sindhi kept it as \emph{pāṛhaṇu} \textquotedblleft{}to teach\textquotedblright{} (CDIAL 8037).
\end{cousinweb}

\begin{grammarlens}[title={What you've built}]
One more word for this chapter.
\end{grammarlens}

\subsection*{Your turn}
\begin{itemize}
\raggedright
  \item \emph{Say it:} \emph{pāṭhayati}
  \item \emph{Say it:} \emph{pāṭhayati}, once more
  \item \emph{Say it:} say \emph{pāṭhayati}
  \item \emph{Recall:} say the Sanskrit for searches, then the Sanskrit for learns, then say \emph{pāṭhayati} again
\end{itemize}

\begin{tcolorbox}[breakable,colback=teal!4,colframe=teal!35!black,title={Before you move on}]
What does \sk{पाठयति} mean? (\textquotedblleft{}Teaches\textquotedblright{}.) Say it once more.
\end{tcolorbox}
\section[\texorpdfstring{\sk{विस्मरति} (vismarati)}{विस्मरति (vismarati)}]{\sk{विस्मरति} (vismarati) — forgets}

\label{lesson:SA-C135-vismarati}



\begin{quote}
Before the new one: say the Sanskrit for learns, then the Sanskrit for teaches.
\end{quote}

\subsection*{You'll want to know: \sk{विस्मरति}}
\textbf{\sk{विस्मरति}} — \emph{vismarati} — \textquotedblleft{}forgets\textquotedblright{}.

\begin{cousinweb}
\emph{vismarati} comes from \emph{vi} \textquotedblleft{}away, apart\textquotedblright{} and the root √smṛ \textquotedblleft{}remember\textquotedblright{} (CDIAL 12021): \textquotedblleft{}remembers away\textquotedblright{}, forgets. MW compares √smṛ with Latin \emph{memor} \textquotedblleft{}mindful\textquotedblright{}, the source of English \emph{memory}; it lives on as Hindi \emph{bisarnā} \textquotedblleft{}to forget\textquotedblright{} (CDIAL 12021).
\end{cousinweb}

\begin{grammarlens}[title={What you've built}]
One more word for this chapter.
\end{grammarlens}

\subsection*{Your turn}
\begin{itemize}
\raggedright
  \item \emph{Say it:} \emph{vismarati}
  \item \emph{Say it:} \emph{vismarati}, once more
  \item \emph{Say it:} say \emph{vismarati}
  \item \emph{Recall:} say the Sanskrit for learns, then the Sanskrit for teaches, then say \emph{vismarati} again
\end{itemize}

\begin{tcolorbox}[breakable,colback=teal!4,colframe=teal!35!black,title={Before you move on}]
What does \sk{विस्मरति} mean? (\textquotedblleft{}Forgets\textquotedblright{}.) Say it once more.
\end{tcolorbox}
\section[\texorpdfstring{\sk{आरोहति} (ārohati)}{आरोहति (ārohati)}]{\sk{आरोहति} (ārohati) — climbs}

\label{lesson:SA-C135-arohati}



\begin{quote}
Before the new one: say the Sanskrit for teaches, then the Sanskrit for forgets.
\end{quote}

\subsection*{You'll want to know: \sk{आरोहति}}
\textbf{\sk{आरोहति}} — \emph{ārohati} — \textquotedblleft{}climbs\textquotedblright{}.

\begin{cousinweb}
\emph{ārohati} comes from \emph{ā} \textquotedblleft{}onto, up to\textquotedblright{} and the root √ruh \textquotedblleft{}climb, rise, grow\textquotedblright{} (MW; CDIAL 1335): \textquotedblleft{}climbs up, mounts\textquotedblright{}. It lives on as Hindi \emph{arohnā} \textquotedblleft{}to mount, ride\textquotedblright{} (CDIAL 1335).
\end{cousinweb}

\begin{grammarlens}[title={What you've built}]
One more word for this chapter.
\end{grammarlens}

\subsection*{Your turn}
\begin{itemize}
\raggedright
  \item \emph{Say it:} \emph{ārohati}
  \item \emph{Say it:} \emph{ārohati}, once more
  \item \emph{Say it:} say \emph{ārohati}
  \item \emph{Recall:} say the Sanskrit for teaches, then the Sanskrit for forgets, then say \emph{ārohati} again
\end{itemize}

\begin{tcolorbox}[breakable,colback=teal!4,colframe=teal!35!black,title={Before you move on}]
What does \sk{आरोहति} mean? (\textquotedblleft{}Climbs\textquotedblright{}.) Say it once more.
\end{tcolorbox}
\section[\texorpdfstring{\sk{अवरोहति} (avarohati)}{अवरोहति (avarohati)}]{\sk{अवरोहति} (avarohati) — climbs down}

\label{lesson:SA-C135-avarohati}



\begin{quote}
Before the new one: say the Sanskrit for forgets, then the Sanskrit for climbs.
\end{quote}

\subsection*{You'll want to know: \sk{अवरोहति}}
\textbf{\sk{अवरोहति}} — \emph{avarohati} — \textquotedblleft{}climbs down\textquotedblright{}.

\begin{cousinweb}
\emph{avarohati} comes from \emph{ava} \textquotedblleft{}down\textquotedblright{} and the root √ruh \textquotedblleft{}climb, rise\textquotedblright{} (MW; CDIAL 817): \textquotedblleft{}climbs down, descends\textquotedblright{}, the mirror of \emph{ārohati}. Pali wore it down to \emph{orohati} \textquotedblleft{}descends\textquotedblright{} (CDIAL 817).
\end{cousinweb}

\begin{grammarlens}[title={What you've built}]
One more word for this chapter.
\end{grammarlens}

\subsection*{Your turn}
\begin{itemize}
\raggedright
  \item \emph{Say it:} \emph{avarohati}
  \item \emph{Say it:} \emph{avarohati}, once more
  \item \emph{Say it:} say \emph{avarohati}
  \item \emph{Recall:} say the Sanskrit for forgets, then the Sanskrit for climbs, then say \emph{avarohati} again
\end{itemize}

\begin{tcolorbox}[breakable,colback=teal!4,colframe=teal!35!black,title={Before you move on}]
What does \sk{अवरोहति} mean? (\textquotedblleft{}Climbs down\textquotedblright{}.) Say it once more.
\end{tcolorbox}
